\documentclass[11pt,a4paper]{article}
\usepackage[utf8]{inputenc}
\usepackage[T1]{fontenc}
\usepackage{geometry}
\geometry{margin=0.75in}
\usepackage{amsmath,amssymb,amsfonts}
\usepackage{booktabs}
\usepackage{hyperref}
\usepackage{enumitem}
\usepackage{xcolor}

% --- Custom Colors ---
\definecolor{primary}{RGB}{0, 51, 102}
\definecolor{secondary}{RGB}{102, 102, 102}

% --- Headers ---
\usepackage{fancyhdr}
\pagestyle{fancy}
\fancyhf{}
\lhead{\textbf{Meeting Handout: Thesis Status Update}}
\rhead{April 7, 2026 | 11:00 AM}
\lfoot{Brhanu Atsbaha}
\cfoot{\thepage}

\begin{document}

\begin{center}
    {\large \textbf{Multi agent systems for image analysis: object and complex scene detection}} \\
    \vspace{0.1cm}
    \textit{Reducing Hallucinations in Historical Illustrations via Cross-Modal Agreement}
\end{center}

\vspace{0.3cm}

\section*{1. Multi-Agent Framing \& Conceptual Goal}
The project has evolved into a \textbf{Multi-Agent System} where the existing object detection pipeline serves as a core agent, complemented by other modal "validators" (CLIP Scene Classification, LLaVA-OneVision VQA, and Kosmos-2.5 Layout Analysis).
\begin{itemize}[noitemsep]
    \item \textbf{Scenario:} High visual noise in archival scans triggers "hallucinations" in foundational VLMs.
    \item \textbf{System Fix:} Require strict \textit{agreement} between independent signals (spatial, semantic, textural) before accepting a detection as realistic. This suppresses false positives and ensures scene-level consistency.
\end{itemize}

\section*{2. Research Questions (RQs)}
As requested by Shereen, current focus is on the following four inquiries:
\begin{enumerate}[noitemsep]
    \item \textbf{RQ1: Domain Adaptation}: How can generative restoration tiers normalize archival noise to enable modern foundational feature extraction?
    \item \textbf{RQ2: Multi-Agent Distillation}: Can an ensemble of "heavy" VLMs (LLaVA-OneVision, Florence-2) be effectively distilled into a real-time YOLO agent for large-scale indexing?
    \item \textbf{RQ3: Semantic Completeness}: To what extent does the fusion of bounding boxes, scene graphs, and OCR provide a verifiable "visual index" for archival researchers?
    \item \textbf{RQ4: Comparative Verification}: How does automated multi-agent consensus perform compared to expert human archival baselines?
\end{enumerate}

\section*{3. Evaluation Infrastructure}
\begin{itemize}[noitemsep]
    \item \textbf{Scene-Aware Agreement (SAA)}: Our primary metric fusing IoUoverlap with CLIP-derived visual and semantic similarity scores.
    \item \textbf{Human Baseline Audit (Tier 3)}: Current next step. A stratified sample (2,000 images) categorized by "VQA object density" to be manually audited for gold-standard validation.
    \item \textbf{Global Mean Agreement (GMA)}: Recent v8 update achieved a GMA of \textbf{0.4704} across 37 discovered classes.
\end{itemize}

\section*{4. Technical Milestone Summary}
\begin{table}[h!]
\centering
\small
\begin{tabular}{lp{10cm}}
\toprule
\textbf{Component} & \textbf{Current Status / Impact} \\
\midrule
Restoration Tier & CycleGAN + Real-ESRGAN + SDXL Refiner completed for full benchmark. \\
VLM Ensemble & Florence-2 and GroundingDINO integrated; achieved \textbf{+11.4\%} consensus boost. \\
YOLO Distillation & Final Student trained on 8,212 consensus-verified images (mAP50: \textbf{0.989}). \\
Taxonomy & Transitioned to 37 "AI-Discovered" Open-Vocab classes to remove bias. \\
\bottomrule
\end{tabular}
\end{table}

\vspace{0.3cm}
\noindent \textbf{Topic for Discussion:} Finalizing the sampling logic for the Tier 3 Human Audit—balancing "High Complexity" (dense 19th-century crowds) vs "Empty Archives" (landscapes).

\end{document}
